% ---------------------------------------------------------------------------
% The switching-cost version's introduction. Poses the proposal's question, says what
% the measurement did to it, and promises what the eight sections after it deliver ---
% in the order the work actually produced them, which puts coverage before the finding
% that supersedes it. Numbers are quoted from the complete paper and carry the same
% registry keys.
% ---------------------------------------------------------------------------

\section{Introduction}
\label{s2:sec:intro}

An observability bill is one of the few software costs that can grow faster than the
system it measures, and the published cases are startling: one company's monitoring
invoice reached \$65M a year, run up, by the account its own engineers give, because
nobody was watching infrastructure spend~\citep{orosz2023}. The obvious response is to
move to something cheaper, and for metrics the cheap option is well established ---
Prometheus is free, widely deployed, and its query language is a de facto standard. Teams
nonetheless stay, and when they explain why, the reason is rarely the data. It is the
dashboards.

A mature observability estate is thousands of saved queries: panels, alerting rules,
service-level objectives, runbook links. Moving vendors means rewriting all of them in
another language, by hand, without a way to tell whether a rewrite preserved what the
original meant. That is a switching cost, and the question worth asking is where it comes
from. If it is incidental --- different keywords, different punctuation --- a translator
removes it and the lock-in was never real. If it follows from how the two languages are
\emph{designed}, no amount of engineering removes it, and the query language is a retention
mechanism whether or not anyone intended it as one.

The question is tractable because query languages are formal languages. Rich treats them
as such (§48.1), and the book's classification --- regular, context-free, decidable --- is
the frame this paper works in~\citep{rich2008}. If the two languages are context-free they
have grammars; if their grammars are unambiguous each query has one parse tree; and where a
construct of one has no meaning-preserving counterpart in the other, the obstruction can be
exhibited as a theorem instead of reported as a limitation.

\paragraph{What was proposed, and what the measurement did to it.}
This work was proposed as a single question --- is the lock-in a syntactic accident or a
semantic necessity? --- with a plan attached: formalise both languages as
$(\text{grammar},\text{semantics})$ pairs, define what it means for a translation between
them to be sound, build a syntax-directed transpiler through a shared intermediate
representation, and measure on real dashboards how much of each language actually
translates. That plan was carried out as written, and three things happened to it.

The first is a correction the proposal could not have anticipated, and it changed the
definition of the central term. Soundness cannot be an equation between two queries
evaluated on one database, because the two systems do not receive the same database: a
Prometheus counter is a running total and a Datadog counter a sequence of per-interval
deltas. Soundness here is therefore stated relative to an \emph{ingestion mapping}
$\phi$, and every theorem names the class of databases it holds on
(\S\ref{s2:sec:framework}).

The second is that the expected answer did not arrive. The proposal took for granted
that some constructs would be semantically inexpressible --- the
\emph{language-design conjecture} --- and the paper reserved a column for them. In the forward direction that column emptied of everything we can
demonstrate, bar $14$ occurrences left undecided, and most of what emptied it were
corrections to our own earlier claims rather than to the vendor's behaviour
(\S\ref{s2:sec:coverage}). The measurement the plan called for was made, and it answers
the question in the negative: on this corpus, in this direction, the switching cost is
work somebody has not done.

The third is that the plan's instruments found something they were not built to find, and
it is the finding this paper is named for. Measuring how much of each language translates
required recording, for every translation, the context the query did not supply. Those
records were kept as qualifications on the coverage number, one per code, and read
together they are a result: a dashboard query does not determine its own meaning. That is
why this paper reports coverage first and the finding second, which is the order the work
produced them in and the reverse of their importance.

\paragraph{What we found.}
Going from Datadog metric queries to PromQL, $96.6\%$ of query occurrences translate. The
shortfall is five named gap codes --- $3.3\%$ carrying a code for a construct this
translator has not implemented, $0.1\%$ that does not parse --- so the ceiling is $99.9\%$
once the first is written, or $99.7\%$ if the $14$ occurrences this paper cannot decide
turn out to be beyond reach after all. Those figures are measured on the corpus the grammar
and the implementation order were fitted to, which \S\ref{s2:sec:threats} states as the
threat it is; frozen and run on a disjoint repository of dashboards by other authors, the
translator reaches $96.7\%$ of $3{,}065$ occurrences.

Translating is not the same as translating correctly, and the distance between the two is
measurable here. Every expression we emit assumes a metric type the query does not state,
and against the vendor's own declarations that inference is wrong for $6.5\%$ of the cases
it can settle --- $491$ occurrences, against a column of demonstrated inexpressibility
holding nothing. The rollup interval is the second and larger case: a Datadog query has no
slot for it, so the interval arrives from the dashboard's time range, and $97.0\%$ of our
translations are of queries that leave it there. PromQL does have the slot and $55.6\%$ of
this corpus's range selectors fill it with a dashboard variable anyway, so the contrast
between the two languages is one of degree and what they share is that the value comes from
a layer outside either of them. Together with the template-variable shares --- $88.6\%$ of
Datadog and $82.5\%$ of PromQL occurrences --- that makes the object a migration ports a
query \emph{schema} plus a resolution layer belonging to the dashboard tool, which is where
the switching cost we measured sits and where neither $\kappa$ nor a gap catalogue can see
it. \S\ref{s2:sec:unsaid} is that case in full.

Coming back is a different matter, and it is where language design does bite. A Datadog
tag filter is a glob and a PromQL tag filter is a regular expression, and globs are a
proper subclass in a way that admits proof rather than measurement.
\S\ref{s2:sec:reverse} states the results and the decision procedure behind them, with the
proofs left to the complete paper. Put beside the forward direction they give an asymmetry:
leaving costs work, returning costs expressiveness.

\paragraph{Scope.}
We study \emph{forward portability}: once a team dual-writes or switches, a translated
query must mean the same thing as the original on the data each system receives from then
on. Migrating stored history is out of scope, and we measure the query layer only --- a
real migration also re-instruments services, re-models cardinality and cost, re-establishes
retention and rebuilds alerting, none of which is measured here. The main direction is
Datadog metric queries into PromQL, because that is the direction teams actually take.
LogQL and Splunk SPL were named in the proposal as stretch targets and remain stretch
targets; no grammar or translator for either appears here.

\paragraph{What this version carries.}
This is the short version of a longer paper, organised around the measurement rather than
around the theorems. \S\ref{s2:sec:background} places the question among the five
literatures it belongs to; \S\ref{s2:sec:framework} fixes the framework;
\S\ref{s2:sec:translator} reports the corpus, the grammars and the translator;
\S\ref{s2:sec:coverage} is the coverage measurement and the gap catalogue;
\S\ref{s2:sec:unsaid} is the paper's result; \S\ref{s2:sec:reverse} states the reverse
direction; and \S\ref{s2:sec:threats} and \S\ref{s2:sec:conclusion} carry the threats and
the answer to the proposal's question. The proofs, the ingestion mapping in full, the
grammar conflict analysis, the differential rig and the live-account experiments are in the
complete paper and cited here by name rather than reproduced.
